\documentclass[12pt]{article}
\usepackage{amsmath, amssymb}
\usepackage{geometry}
\geometry{margin=1in}

\title{CSE400: Fundamentals of Probability in Computing \ Joint Random Variables and Distributions}
\author{}
\date{}

\begin{document}

\maketitle

\section{Motivation}

In many applications, we study two or more random variables simultaneously.

\textbf{Examples:}
\begin{itemize}
\item Height and weight
\item IQ and birthweight
\item Exercise frequency and heart disease
\item Signal-to-noise ratio (SNR) and throughput
\item Rainfall and temperature
\end{itemize}

These variables may be \textbf{dependent} or \textbf{independent}. Their relationship is described using a \textbf{joint distribution}.

\textbf{Key idea:} Joint distributions allow us to compute probabilities involving multiple variables and analyze relationships using covariance and correlation.

\section{Joint PMF (Discrete Case)}

Let:
[
X \in {x_1, x_2, \dots, x_n}, \quad Y \in {y_1, y_2, \dots, y_m}
]

\textbf{Definition:}
[
p(x_i, y_j) = P(X = x_i, Y = y_j)
]

\textbf{Properties:}
\begin{align*}
0 \le p(x_i, y_j) \le 1 \
\sum_{i=1}^{n} \sum_{j=1}^{m} p(x_i, y_j) = 1
\end{align*}

\section{Example: Two Dice}

Let:
[
X = \text{first die}, \quad T = \text{sum of both dice}
]

Each outcome has probability:
[
\frac{1}{36}
]

\textbf{Observation:} $X$ and $T$ are \textbf{not independent} since $T$ depends on $X$.

\section{Joint PDF (Continuous Case)}

Let:
[
X \in [a,b], \quad Y \in [c,d]
]

\textbf{Definition:} A function $f(x,y)$ is a joint PDF if:
[
P((X,Y) \text{ in small region}) \approx f(x,y),dx,dy
]

\textbf{Properties:}
\begin{align*}
f(x,y) \ge 0 \
\int_c^d \int_a^b f(x,y),dx,dy = 1
\end{align*}

\textbf{Note:} $f(x,y)$ may be greater than 1 (it is a density).

\section{Worked Example}

Given:
[
f(x,y) = 4xy, \quad 0 \le x \le 1,; 0 \le y \le 1
]

\subsection{Check Validity}

\begin{align*}
\int_0^1 \int_0^1 4xy,dx,dy
&= \int_0^1 \left( \int_0^1 4xy,dx \right) dy \
&= \int_0^1 2y,dy \
&= \left[y^2\right]_0^1 = 1
\end{align*}

Hence, valid joint PDF.

\subsection{Compute Probability}

Event:
[
A = {X < 0.5,; Y > 0.5}
]

\begin{align*}
P(A)
&= \int_{0.5}^{1} \int_{0}^{0.5} 4xy,dx,dy \
&= \int_{0.5}^{1} \left[ 2x^2 y \right]*{0}^{0.5} dy \
&= \int*{0.5}^{1} 2y(0.25),dy \
&= \int_{0.5}^{1} \frac{y}{2},dy \
&= \left[\frac{y^2}{4}\right]_{0.5}^{1} \
&= \frac{1}{4} - \frac{0.25}{4} = \frac{3}{16}
\end{align*}

\section{Joint CDF}

\textbf{Definition:}
[
F(x,y) = P(X \le x, Y \le y)
]

\subsection{Continuous Case}

[
F(x,y) = \int_c^y \int_a^x f(u,v),du,dv
]

\textbf{Recover PDF:}
[
f(x,y) = \frac{\partial^2}{\partial x \partial y} F(x,y)
]

\subsection{Discrete Case}

[
F(x,y) = \sum_{x_i \le x} \sum_{y_j \le y} p(x_i, y_j)
]

\section{Properties of Joint CDF}

\begin{itemize}
\item Non-decreasing in both $x$ and $y$
\item $\lim_{x,y \to -\infty} F(x,y) = 0$
\item $\lim_{x,y \to \infty} F(x,y) = 1$
\end{itemize}

\end{document}
